%% ASPLOS 2026 submission template
%% Class:  acmart (sigplan proceedings)
%% Per ASPLOS'26 CFP: \documentclass[sigplan,nonacm]{acmart}
%% Body text 10pt. Page numbers enabled via \settopmatter{printfolios=true}.
%% Bibliography: ACM-Reference-Format.
%% This is the Project Checkpoint document (2--3 pages).

\documentclass[sigplan,nonacm,review]{acmart}

\settopmatter{printfolios=true,printccs=false,printacmref=false}

%% ---- Packages -----------------------------------------------------------
\usepackage{booktabs}
\usepackage{subcaption}
\usepackage{balance}

%% ---- Metadata (stripped by nonacm, but kept for completeness) -----------
\acmConference[ASPLOS '26]{31st International Conference on Architectural
  Support for Programming Languages and Operating Systems}{March 22--26,
  2026}{Pittsburgh, PA, USA}

%% ---- Title --------------------------------------------------------------
\title{Optimizing Qwen2.5-Coder Inference on Tenstorrent Blackhole:\\
       Project Checkpoint}

\author{Rassul Magauin}
\author{Assylzhan Khamiyev}
\author{Mohammed Rashed Ali Yahmoor Alshehhi}
\author{Sultan Akimaliyev}
\affiliation{%
  \institution{Mohamed bin Zayed University of Artificial Intelligence}
  \city{Abu Dhabi}\country{UAE}}

\renewcommand{\shortauthors}{Magauin et al.}

%% ---- Abstract -----------------------------------------------------------
\begin{abstract}
We report the first checkpoint of an inference optimization study of
Qwen2.5-Coder on the Tenstorrent Blackhole~p150 accelerator using the
\texttt{tt-metal} SDK. Starting from the in-tree \texttt{tt\_transformers}
port with the stock Blackhole tuning recipe, we establish a device-level
decode baseline on a four-device tensor-parallel mesh, identify the
dominant kernels via the Tracy device profiler and the
\texttt{tt-perf-report} utility recommended by the course staff, and rank
three concrete optimization targets by projected return on engineering
effort. Our measurements reveal that (i)~the sampling-stage TopK operator
accounts for 45--60\% of per-token device kernel time but executes on a
single Tensix core out of 130---an under-utilization artifact rather than
a model bottleneck; (ii)~the MLP feed-forward block is the largest
\emph{real} bottleneck at $\sim$31\% of device time, dominated by the
gate/up and down projections; and (iii)~the fused scaled-dot-product
attention kernel is negligible ($<$1.5\%), validating that existing
flash-attention-style fusions on Blackhole already perform as intended.
We propose parallelizing TopK, switching the MLP matmuls from HiFi4 to
HiFi2 math fidelity, and fusing the gate/up projections as our next
steps.
\end{abstract}

\begin{document}
\maketitle

%% =========================================================================
\section{Introduction and Goals}
%% =========================================================================
Large language models tailored to code---such as Qwen2.5-Coder~\cite{qwen25coder}---are
increasingly deployed on non-GPU accelerators. Tenstorrent's Blackhole
p150 is an open-architecture, RISC-V-based tensor processor whose 130
worker Tensix cores each expose a matrix unit (FPU), a vector unit
(SFPU), five Baby RISC-V CPUs, and 1.5~MB of L1 SRAM connected by a
2D NoC~\cite{ttmetal}. Unlike GPU stacks, \texttt{tt-metal} exposes
kernels, circular buffers, and NoC scheduling directly to the developer,
which creates both an optimization opportunity and a steep perf-debugging
curve.

\paragraph{Project direction.}
Following the teaching staff's redirect, we pursue \emph{Direction~B:
Optimization}: we take the already-ported \texttt{tt\_transformers}
implementation of Qwen2.5 as a starting point and measure, rank, and
remove bottlenecks specific to Blackhole. A from-scratch port was
rejected because \texttt{tt\_transformers} already supports the Qwen2.5
family.

\paragraph{Contributions of this checkpoint.}
\begin{enumerate}
  \item A reproducible device-level profiling pipeline for
        \texttt{tt\_transformers} decode on a four-device Blackhole mesh,
        combining the Tracy profiler with the \texttt{tt-perf-report}
        utility recommended by TA Dongning Ma.
  \item A per-operator breakdown of Qwen2.5-Coder-7B-Instruct decode
        against the stock tuned recipe, distinguishing
        under-utilization artifacts (TopK) from real model bottlenecks
        (MLP).
  \item A ranked list of three concrete, Blackhole-aware optimization
        targets with projected first-order ROI, to be executed and
        measured before the final report.
\end{enumerate}

%% =========================================================================
\section{Background}
%% =========================================================================
\paragraph{Blackhole p150.}
Blackhole exposes an $8{\times}10$ compute grid (130 worker cores
visible to \texttt{tt-metal}, compared to Wormhole's $8{\times}8$),
1.5~MB of L1 per core, and eight DRAM banks. The in-tree
\texttt{tt\_transformers} framework gates several Blackhole-specific
paths via \texttt{is\_blackhole()}, most notably \texttt{prefill\_len\_cutoff}
$=\!512$ (vs.\ 1024 on Wormhole) and a Blackhole-only DRAM weight
prefetcher currently wired for Llama-3.1-8B.

\paragraph{Qwen2.5-Coder-7B-Instruct.}
Qwen2.5-Coder-7B shares its base transformer with
Qwen2.5-7B-Instruct: 28 layers, hidden size 3584, MLP intermediate
18944, 28 query heads with 4 grouped key/value heads, and
\texttt{rope\_theta} $=10^{6}$. This architectural parity lets us use
the \texttt{performance\_decoder\_config.json} tuning recipe that
\texttt{tt-metal} already ships for Qwen2.5-7B.

\paragraph{Fusion model of \texttt{tt-metal}.}
\texttt{tt-metal} fuses via circular-buffer choreography rather than
monolithic kernels. A \emph{reader} kernel streams tiles from DRAM or a
neighbor core's L1 into L1 circular buffers; a \emph{compute} kernel
drains them through the FPU/SFPU; and a \emph{writer} kernel streams
results out. Fusion is then a property of the CB topology, not a single
file. We rely on five canonical fusion patterns already present in
tt-metal: fused SDPA, matmul$+$bias$+$activation, distributed RMSNorm,
async AllGather$+$matmul, and paged KV-cache update.

%% =========================================================================
\section{Methodology}
%% =========================================================================
\paragraph{Hardware.} One machine with four Blackhole p150 devices at
\texttt{/dev/tenstorrent/\{0..3\}}, running \texttt{tt-metal}
built with \texttt{ENABLE\_TRACY=ON}.

\paragraph{Workload.} We drive the canonical
\texttt{models/tt\_transformers/demo/simple\_text\_demo.py} benchmark:
decode mode, batch~1, 1024-token context, 10 transformer layers, paged
attention on, and the stock
\texttt{performance\_decoder\_config.json}. The \texttt{blackhole-4}
pytest fixture auto-selects a 4-way tensor-parallel mesh. Weights are
local \texttt{bf16} safetensors downloaded from Hugging Face.

\paragraph{Profiling pipeline.} We wrap the pytest invocation with
\texttt{python -m tracy -p -r -a device\_kernel\_duration}, which emits
a per-op CSV under
\texttt{generated/profiler/reports/.../ops\_perf\_results\_*.csv}
(4{,}916 rows across two trace-replay sessions, compile phase
excluded). We post-process the CSV with two complementary tools: a
small pandas script for per-op aggregates, and the
\texttt{tt-perf-report}~v1.2.2 utility distributed as a separate
repository~\cite{ttperfreport} (explicitly requested by our TA) for
stacked per-op-class breakdowns and math-fidelity analysis.

%% =========================================================================
\section{Baseline Results}
%% =========================================================================
Table~\ref{tab:baseline} reports the per-iteration device-kernel-time
breakdown for Qwen2.5-Coder-7B-Instruct decode on a four-way Blackhole
tensor-parallel mesh (device~0, steady-state trace replay, compile
phase excluded). Total per-decode-iteration device kernel time is
\textbf{12.1~ms} on the 7B-Instruct run and 10.4~ms on the Coder-7B
run; the absolute delta is attributable to compile-cache warm-up
between back-to-back runs (op counts and shapes are bit-identical, and
the TopK kernel time matches to the nanosecond), so the
\emph{qualitative} ranking is robust across runs.

\begin{table}[t]
\centering
\small
\caption{Per-iteration device-0 kernel time, Qwen2.5-Coder-7B-Instruct
  decode, 4-way Blackhole TP mesh, 10 layers, stock tuned recipe.
  TopK is a sampling-stage op, not a transformer-layer op.}
\label{tab:baseline}
\begin{tabular}{lrr}
\toprule
Operator / group & Time (ms) & \% of total \\
\midrule
TopK (sampling, 1-core)          & 5.48 & 45.3 \\
MLP gate$+$up projections        & 2.41 & 19.9 \\
MLP down projection              & 1.40 & 11.5 \\
Attention QKV projections        & 1.03 &  8.5 \\
CCL (AllGather $+$ ReduceScatter)& 0.69 &  5.7 \\
Layout/tiling/eltwise glue       & 0.46 &  3.8 \\
Attention output projection      & 0.28 &  2.3 \\
Fused SDPA decode $+$ RoPE       & 0.15 &  1.3 \\
RMSNorm                          & 0.13 &  1.1 \\
Other                            & 0.08 &  0.6 \\
\midrule
\textbf{Total per iter.}         & \textbf{12.11} & \textbf{100.0} \\
\bottomrule
\end{tabular}
\end{table}

%% =========================================================================
\section{Bottleneck Analysis}
%% =========================================================================
\paragraph{TopK is an under-utilization artifact, not a layer bug.}
\texttt{TopKDeviceOperation} consumes 45--60\% of per-token device
kernel time across both profiled runs, with a shape of
$[1,1,32,38016]$ and $k{=}32$. The kernel runs on a \emph{single}
Tensix core out of the 130 available on Blackhole. This is not a
transformer-layer bottleneck---it is a sampling-stage mapping bug that
an existing multi-core reduction template should be able to remove
nearly in full. We treat it as the highest-ROI, lowest-risk item on
the list.

\paragraph{The MLP block is the real hot path.}
Excluding TopK, the MLP feed-forward block accounts for
$\sim$31\% of device kernel time ($\approx$3.81~ms/iter), split
between the gate/up projections ($3584 \!\to\! 5120$, 20 calls/iter)
and the down projection ($5120 \!\to\! 3584$). The single slowest
individual op is the down projection at 139.8~$\mu$s
max. Per the tt-metal model-config database, these matmuls currently
run in HiFi4 math fidelity on BF16/BFP8 operands.

\paragraph{Attention is already fused and cheap.}
The fused SDPA-decode kernel, head reshapes, RoPE, and KV-cache update
together contribute only 1.3\% of per-iter kernel time. This confirms
that the flash-attention-style CB choreography in
\texttt{ttnn/cpp/ttnn/operations/transformer/sdpa/} is already doing
what we would hope; attention is explicitly \emph{not} a target.

\paragraph{Tensor-parallel overhead is modest.}
AllGather plus ReduceScatter (the TP tax) accounts for 5.7\%. This is
headroom worth chasing only after the MLP wins are realized.

%% =========================================================================
\section{Planned Optimizations and Projected ROI}
%% =========================================================================
We rank three candidates by (a)~projected kernel-time recovery and
(b)~engineering effort.

\begin{enumerate}
\item \textbf{Parallelize / fuse TopK.} Move \texttt{TopKDeviceOperation}
      off a single core onto a multi-core reduction template, or fuse
      it directly into \texttt{SamplingDeviceOperation}. Projected
      recovery: most of the 45--60\% TopK slice. Effort:
      \emph{low}---this is a dispatch-and-shard problem, not a new
      kernel. (Per-layer scaling: frees time regardless of layer
      count, since sampling is per-token.)
\item \textbf{HiFi4 $\to$ HiFi2 for MLP matmuls.} The Blackhole FPU
      exposes a HiFi2 math-fidelity mode that trades the lowest
      activation bit for roughly 2$\times$ matmul throughput. The
      \texttt{tt-perf-report} tool flagged this as the highest-impact
      single-line change available against the stock recipe.
      Projected recovery: up to half of the $\sim$3.8~ms MLP slice
      ($\approx$1.9~ms/iter, $\approx$15\% of TopK-excluded kernel
      time). Effort: \emph{very low}---one field in
      \texttt{performance\_decoder\_config.json}. Risk: must verify
      with \texttt{ci-token-matching} that end-to-end generation
      quality is preserved.
\item \textbf{Fuse MLP gate$+$up projections.} Merge the two
      $3584\!\to\!5120$ projections into one $3584\!\to\!10240$ matmul,
      halving dispatch overhead and activation round-trips. Projected
      recovery: the second-largest slice after HiFi2. Effort:
      \emph{medium}---requires checking whether the stock config
      already enables this and, if not, a model-code edit plus a
      perf/accuracy re-run.
\end{enumerate}

%% =========================================================================
\section{Roadblocks and Risks}
%% =========================================================================
\paragraph{Hardware monitoring.} \texttt{tt-smi} is not installed on our
box; we confirm device health via first-program-launch runtime logs.
Installing \texttt{tt-smi} is a next step.

\paragraph{Compile-cache variance.} Back-to-back runs of architecturally
identical models showed a $\sim$44\% per-op difference in Matmul mean
time, attributable to compile-cache state. We will rerun each
optimization experiment from a clean cache and report matched-cache
deltas.

\paragraph{Accuracy regression from HiFi2.} HiFi2 drops the lowest
activation mantissa bit. We will gate the change behind the in-tree
\texttt{ci-token-matching} parameterization and report token-match
rates against the HiFi4 reference.

\paragraph{Coder-specific validation.} Since Qwen2.5-Coder-7B is
architecturally identical to Qwen2.5-7B-Instruct, our ``Coder-ness''
must be demonstrated at the accuracy/output level (code-generation
prompts), not at the kernel level. We will include a small
code-completion token-match check in the final report.

%% =========================================================================
\section{Next Steps}
%% =========================================================================
Between now and the final report we will:
(1)~land the TopK multi-core fix and re-profile;
(2)~measure HiFi2 MLP speedup and verify token-match rate;
(3)~evaluate MLP gate$+$up fusion;
(4)~characterize single-device (non-TP) decode to isolate CCL
overhead; and
(5)~profile prefill in addition to decode, since the Blackhole
prefill-length cutoff of 512 tokens may expose a distinct set of
hot paths.

%% =========================================================================
\balance
\bibliographystyle{ACM-Reference-Format}
\bibliography{refs}

\end{document}
